% Original 2001 IEL examples; source-native syntax and separately supplied proofs.
\section{Evocation and implication patterns}
\label{sec:p18-erotetic-patterns-and-composition}

A sound question need not be evoked by the current information: one of its
answers may already follow. Conversely, a legitimate auxiliary question need
not itself be evoked; both of its responses may lead to the same original
answer. The distinction is visible in the concrete patterns of
Wi\'sniewski's original presentation
\cite[secs.~6,8--9.1]{Wisniewski2001QuestionsInferences}.
This section gives their full finite families and first-order extents.
The semantic relations are those of
Definitions~\ref{def:p18-ei-setting} and~\ref{def:p18-ei-relation}.
The original schemata are semantic warrants, not an asserted complete
syntactic calculus for all erotetic implication.

\subsection{The original languages and notation}

\begin{hbdefinition}{def:p18-ep-notation}{Source-native languages and schema variables}
The source's \(L_1\) has the classical propositional declaratives and the
explicit questions \(?\{A_1,\ldots,A_n\}\), \(n>1\), with pairwise
syntactically distinct answers. Its Boolean signs are
\(\neg,\supset,\land,\lor,\equiv\): here \(\supset\) is material
implication and \(\equiv\) is the Boolean biconditional, not set inclusion
or elementary equivalence. Their truth functions are the same as those
specified for \(\to,\leftrightarrow\) in the preceding declarative
presentation, but the source's printed signs are retained here.

Write \(?A\) for \(?\{A,\neg A\}\). The source's binary conjunctive
question is \(?\pm|A,B|\), with direct-answer set
\[
 \{A\land B,A\land\neg B,\neg A\land B,\neg A\land\neg B\}.
\]
The sign \(\pm\) is part of this question abbreviation, not addition.
The expression \(A\otimes B\) in the schemata has
\(\otimes\in\{\land,\lor,\supset,\equiv\}\); this \(\otimes\)
is a metalinguistic connective parameter, not a linear-logic connective.
The source's \(\alpha A\) denotes either \(A\) or \(\neg A\), and
\(\alpha B\) either \(B\) or \(\neg B\). The choice is fixed within
each instance. Here \(A,B,A_i\) are formulas; the lower-case
\(p,q,r\) in the propositional evocation examples are distinct variables.

The source's \(L_2\) has monadic first-order declaratives with equality,
constants, no function symbols, and all nonempty-domain interpretations.
The explicit answers are closed sentences. \(Ax\) has \(x\) as its only
free variable; \(A(x/c)\) is the closed result of substituting constant
\(c\) for its free occurrences. In this notation the slash is a
substitution separator. The existential and general wh-questions are
\(?S(Ax)\) and \(?U(Ax)\), with the answer sets of
Definition~\ref{def:p18-ei-wh-language}. On the metalanguage level
\(S(Ax)\) and \(U(Ax)\) designate those sets; on the object-language
level the same letters occur in the question expressions.

Distinct constant metavariables designate distinct constant symbols, not
necessarily distinct denotations. In the evocation examples \(P,R\) are
distinct predicate symbols. Quantifier-free implication patterns apply
also to closed \(L_2\)-sentences by their classical truth clauses; they
do not depend on predicates being independent Boolean variables.

The source abbreviates \(E(\{C_1,\ldots,C_k\},Q)\) as
\(C_1,\ldots,C_k\vdash_E Q\), and
\(\operatorname{Im}(Q,\{C_1,\ldots,C_k\},R)\) as
\(Q,C_1,\ldots,C_k\vdash_{\operatorname{Im}}R\).
These turnstiles abbreviate semantic relations, not PMCE proof judgments.
\end{hbdefinition}

\begin{hbdefinition}{def:p18-ep-families}{Finite families and complete-list premises}
For an answer list \(A_1,\ldots,A_n\), put
\(J_n=\{1,\ldots,n\}\),
\(D_n=A_1\lor\cdots\lor A_n\), and
\(K_n=A_1\land\cdots\land A_n\), with a fixed displayed bracketing.
For \(1\leq i<n\), put
\(D^-_i=A_1\lor\cdots\lor A_i\) and
\(D^+_i=A_{i+1}\lor\cdots\lor A_n\).
Both blocks are nonempty. All answer lists in a schema must still satisfy
the source's syntactic question conditions, including the lists of negations.

For a finite nonempty list \(\vec c=(c_1,\ldots,c_k)\) of distinct
constant symbols, define the declarative abbreviations
\[
 N_{\vec c}(x)=(x=c_1\lor\cdots\lor x=c_k),\qquad
 H_{\vec c}=\forall x(Px\equiv N_{\vec c}(x)),
\]
\[
 W_{\vec c}=Pc_1\land\cdots\land Pc_k\land
                 \forall x(Px\supset N_{\vec c}(x)).
\]
The last formula is a direct \(U(Px)\)-answer. The finite family in
source (6.25) is made explicit as \(H_{\vec c^{1}}\lor\cdots\lor
H_{\vec c^{m}}\), with \(m\geq2\), each list length greater than
one, and distinct constant symbols across the alternative lists, as in
the source's distinct-metavariable convention. No unique-names axiom is
added. In the sheet the first list is \((c_1,\ldots,c_n)\) and the
last is \((c'_1,\ldots,c'_k)\), as printed in the source.
When referring to a candidate complete-list answer,
\(\vec t\) is any finite nonempty list of distinct constant symbols.
\end{hbdefinition}

\begin{hbdefinition}{def:p18-ep-analytic}{The source's analytic implication condition}
A question \(Q\) analytically implies \(R\) exactly when it purely
implies \(R\) and every immediate subformula of every answer to \(R\)
is either a subformula of an answer to \(Q\), or the negation of such
a subformula. A formula is a subformula of itself.
For the exact classifications below write
\[
 \mathcal H(Q)=\{C:C\text{ is a subformula of some }D\in dQ\}
 \cup\{\neg C:C\text{ is such a subformula}\}.
\]
This is a syntactic set. Boolean equivalence does not add new subformulas.
Immediate subformulas are the children of the outermost connective;
an atomic formula has none
\cite[p.~27, Definition~11 and note~14]{Wisniewski2001QuestionsInferences}.
\end{hbdefinition}

The following sheet retains the source's numbering and all the remaining
evocation and implication schemata. Its finite-list abbreviations have just
been defined; the theorem statements below specify their full scope and
proofs. It is a sheet of source semantic warrants rather than a claim of
derivability by a common primitive calculus.
\calculussheet{sheet:p18-erotetic-patterns}{Original IEL evocation and implication patterns}
 {Source \(L_1,L_2\) and native \(\supset,\equiv,?\pm|A,B|\).
 These are semantic E/Im warrants, not a complete primitive calculus.
 Every resulting question has at least two distinct closed answers.
 Families and notation: Definitions~\ref{def:p18-ep-notation}--\ref{def:p18-ep-families}.
 Source numbers are retained; proofs are supplied in the matched pattern annex.}

\hbrule{rule:p18-ep-e05}{E (6.5)}{p\lor\neg p}{?p}
 {Distinct propositional variables in E (6.5)--(6.15).
  All eleven proofs: Annex~\ref{annex:p18.ep.propositional-evocation}.}
\hbrule{rule:p18-ep-e06}{E (6.6)}{p\lor q}{?p}{\none}
\hbrule{rule:p18-ep-e07}{E (6.7)}{p\lor q}{?q}{\none}
\hbrule{rule:p18-ep-e08}{E (6.8)}{p\supset q}{?p}{\none}
\hbrule{rule:p18-ep-e09}{E (6.9)}{p\supset q}{?q}{\none}
\hbrule{rule:p18-ep-e10}{E (6.10)}{p\equiv q}{?\{p\land q,\neg p\land\neg q\}}{\none}
\hbrule{rule:p18-ep-e11}{E (6.11)}{p\lor q}{?\{p,q\}}{\none}
\hbrule{rule:p18-ep-e12}{E (6.12)}{p\lor q}
 {?\{p\land q,p\land\neg q,\neg p\land q\}}{\none}
\hbrule{rule:p18-ep-e13}{E (6.13)}{p\supset(q\lor r)\qquad p}{?\{q,r\}}{\none}
\hbrule{rule:p18-ep-e14}{E (6.14)}{(p\land q)\supset r\qquad\neg r}{?\{\neg p,\neg q\}}{\none}
\hbrule{rule:p18-ep-e15}{E (6.15)}{(p\lor q)\lor r}{?\{p,q\lor r\}}{\none}

\hbrule{rule:p18-ep-e16}{E (6.16)}{Pc_1,\ldots,Pc_n}{?\forall xPx}
 {Finite \(n\geq1\); all nonempty-domain \(L_2\)-models, not domain closure.
  Annex~\ref{annex:p18.ep.firstorder-polar}.}
\hbrule{rule:p18-ep-e17}{E (6.17)}{Pc_1\land Rc_1,\ldots,Pc_n\land Rc_n}
 {?\forall x(Px\supset Rx)}
 {Finite \(n\geq1\), distinct predicates. Annex~\ref{annex:p18.ep.firstorder-polar}.}
\hbrule{rule:p18-ep-e18}{E (6.18)}{Pc_1\land Rc_1,\ldots,Pc_n\land Rc_n,\ Pc_{n+1}}
 {?Rc_{n+1}}
 {Distinct constant symbols and predicates, no distinct-denotation axiom.
  Annex~\ref{annex:p18.ep.firstorder-polar}.}
\hbrule{rule:p18-ep-e19}{E (6.19)}{Pc_1,Pc_2}{?c_1=c_2}
 {Distinct symbols may denote the same individual. Annex~\ref{annex:p18.ep.firstorder-polar}.}
\hbrule{rule:p18-ep-e20}{E (6.20)}{\exists xPx}{?\{Pc,\exists x(Px\land x\ne c)\}}
 {Named-or-residual witness. Annex~\ref{annex:p18.ep.named-evocation}.}
\hbrule{rule:p18-ep-e21}{E (6.21)}{\exists x(Px\land(x=c_1\lor\cdots\lor x=c_n))}{?\{Pc_1,\ldots,Pc_n\}}
 {\(n>1\), distinct constant symbols. Annex~\ref{annex:p18.ep.named-evocation}.}
\hbrule{rule:p18-ep-e22}{E (6.22)}{\neg\forall x((x=c_1\lor\cdots\lor x=c_n)\supset Px)}
 {?\{\neg Pc_1,\ldots,\neg Pc_n\}}
 {\(n>1\), distinct constant symbols. Annex~\ref{annex:p18.ep.named-evocation}.}
\hbrule{rule:p18-ep-e23}{E (6.23)}{\exists x(Px\land(x=c_1\lor\cdots\lor x=c_n))}{?S(Px)}
 {\(n>1\); proof of non-entailment for every language constant.
  Annex~\ref{annex:p18.ep.named-evocation}.}
\hbrule{rule:p18-ep-e24}{E (6.24)}{\neg\forall x((x=c_1\lor\cdots\lor x=c_n)\supset Px)}{?S(\neg Px)}
 {\(n>1\); the whole wh-answer set, not just the finite negative list.
  Annex~\ref{annex:p18.ep.named-evocation}.}
\hbrule{rule:p18-ep-e25}{E (6.25)}
 {\forall x(Px\equiv(x=c_1\lor\cdots\lor x=c_n))\lor\cdots\lor
  \forall x(Px\equiv(x=c'_1\lor\cdots\lor x=c'_k))}{?U(Px)}
 {\(m\geq2\), \(n,k>1\), all intermediate alternative list lengths
  greater than one; distinct
  symbols across lists as in the source convention. Every candidate-list
  proof: Annex~\ref{annex:p18.ep.complete-list}.}

\hbrule{rule:p18-ep-im01}{Im (8.1)}{?\neg A}{?A}
 {Regular and analytic. Annexes~\ref{annex:p18.ep.pure-patterns},
  \ref{annex:p18.ep.analytic-profile}.}
\hbrule{rule:p18-ep-im02}{Im (8.2)}{?A}{?\neg A}
 {Regular and analytic. Same complete proofs.}
\hbrule{rule:p18-ep-im03}{Im (8.3)}{?\{A,B\}}{?\{B,A\}}
 {Distinct answers; regular and analytic. Same complete proofs.}
\hbrule{rule:p18-ep-im04}{Im (8.4)}{?\{A,B\lor C\}}{?\{A,B,C\}}
 {Arbitrary formulas; both questions well formed. Regular and analytic.
  Annex~\ref{annex:p18.ep.general-refinement}.}
\hbrule{rule:p18-ep-im05}{Im (8.5)}{?\{A,B,C\}}{?\{A,B\lor C\}}
 {Arbitrary formulas; both questions well formed. Analytic, not generally
  regular. Annex~\ref{annex:p18.ep.general-refinement}.}
\hbrule{rule:p18-ep-im07}{Im (8.7)}{?\pm|A,B|}{?(A\otimes B)}
 {\(\otimes\in\{\land,\lor,\supset,\equiv\}\).
  Analytic exactly for conjunction; not generally regular.
  Annexes~\ref{annex:p18.ep.pure-patterns},\ref{annex:p18.ep.analytic-profile}.}
\hbrule{rule:p18-ep-im10}{Im (8.10)}{?A}{?\pm|A,B|}
 {Regular; analytic iff both \(B,\neg B\in\mathcal H(?A)\).
  Same complete proofs; independent distinct atoms give non-analytic instances.}
\hbrule{rule:p18-ep-im11}{Im (8.11)}{?B}{?\pm|A,B|}
 {Regular; analytic iff both \(A,\neg A\in\mathcal H(?B)\).
  Same complete proofs.}
\hbrule{rule:p18-ep-im12}{Im (8.12)}{?(A\otimes B)\qquad\alpha A}{?B}
 {\(\alpha A\in\{A,\neg A\}\), fixed in each instance;
  all four connectives, regular. Annex~\ref{annex:p18.ep.signed-cofactors}.}
\hbrule{rule:p18-ep-im13}{Im (8.13)}{?(A\otimes B)\qquad\alpha B}{?A}
 {\(\alpha B\in\{B,\neg B\}\), fixed; all four connectives, regular.
  Annex~\ref{annex:p18.ep.signed-cofactors}.}
\hbrule{rule:p18-ep-im14}{Im (8.14)}{?\{\neg A,A\land B,A\land\neg B\}}{?A}
 {Empty basis even though printed in the nonempty-basis subsection.
  Not generally regular. Annex~\ref{annex:p18.ep.polar-refinements}.}
\hbrule{rule:p18-ep-im15}{Im (8.15)}{?\{A_1,\ldots,A_n\}\qquad A_1\lor\cdots\lor A_n}{?A_i}
 {\(n>1\), \(i\in J_n\); regular if \(n=2\).
  Annex~\ref{annex:p18.ep.polar-refinements}.}
\hbrule{rule:p18-ep-im16}{Im (8.16)}
 {?\{A_1,\ldots,A_n\}\qquad A_1\lor\cdots\lor A_n\qquad\neg(A_1\land\cdots\land A_n)}
 {?\{\neg A_1,\ldots,\neg A_n\}}
 {\(n>1\); regular if \(n=2\). Each complementary witness and both
  lost-premise cases: Annex~\ref{annex:p18.ep.denial}.}
\hbrule{rule:p18-ep-im18}{Im (8.18)}
 {?\{A_1,\ldots,A_n\}\qquad B\supset(A_1\lor\cdots\lor A_i)\qquad
  \neg B\supset(A_{i+1}\lor\cdots\lor A_n)}{?B}
 {\(1\leq i<n\); regular if \(n=2\). Annex~\ref{annex:p18.ep.partition-patterns}.}
\hbrule{rule:p18-ep-im19}{Im (8.19)}
 {?\{A_1,\ldots,A_n\}\qquad B\supset(A_1\lor\cdots\lor A_i)\qquad
  C\supset(A_{i+1}\lor\cdots\lor A_n)\qquad B\lor C}
 {?\{B,C\}}
 {\(1\leq i<n\); regular if \(n=2\). Same complete block proof.}
\hbrule{rule:p18-ep-im20}{Im (8.20)}
 {?\{A_1,\ldots,A_n\}\qquad(A_1\lor\cdots\lor A_i)\equiv B\qquad
  (A_{i+1}\lor\cdots\lor A_n)\equiv C}{?\{B,C\}}
 {\(1\leq i<n\); regular if \(n=2\). Same complete block proof.}
\hbrule{rule:p18-ep-im21}{Im (8.21)}
 {?\{A_1,\ldots,A_n\}\qquad(A_1\lor\cdots\lor A_i)\supset B\qquad
  (A_{i+1}\lor\cdots\lor A_n)\supset\neg B\qquad A_1\lor\cdots\lor A_n}{?B}
 {\(1\leq i<n\); regular if \(n=2\). The first displayed disjunction
  is the antecedent of its conditional, not an extra premise. Same proof.}
\hbrule{rule:p18-ep-im22}{Im (8.22)}
 {?A(x/c)\qquad\forall x(Ax\equiv(Bx\land Cx))}{?\pm|B(x/c),C(x/c)|}
 {Arbitrary \(Ax,Bx,Cx\) with only x free; closed substitution, no capture;
  regular. Annex~\ref{annex:p18.ep.arbitrary-formulas}.}
\hbrule{rule:p18-ep-im26}{Im (8.26)}
 {?S(Ax)\qquad\exists xAx\qquad\forall x(Ax\supset(x=c_1\lor\cdots\lor x=c_n))}{?A(x/c_i)}
 {\(n>1\), distinct constant symbols, \(i\in J_n\); arbitrary Ax with
  only x free. Positive singleton and negative complementary witnesses;
  regular if \(n=2\). Annex~\ref{annex:p18.ep.wh-polar}.}


\subsection{All propositional evocation examples}

\begin{hbtheorem}{thm:p18-ep-propositional-evocation}{proved=annex:p18.ep.propositional-evocation}{Coverage and non-entailment in (6.5)--(6.15)}
All eleven propositional evocation entries in the sheet hold for distinct
\(p,q,r\), under every Boolean valuation of \(L_1\).
Each supplies coverage and a separate countervaluation for every direct
answer. In (6.12), the answers \(p\land q,p\land\neg q,\neg p\land q\)
include simultaneous truth of \(p,q\), so no exclusive-p-or-q premise
is added. In (6.11), p and q themselves can both be true: evocation
does not require exclusive direct answers. In (6.13) the first premise is
\(p\supset(q\lor r)\), and in (6.14) it is
\((p\land q)\supset r\), with the indicated parentheses.
\end{hbtheorem}

The polar examples (6.5)--(6.9) illustrate that Boolean coverage of
\(?p\) alone is automatic, while evocation still needs witnesses for
both answers not being entailed. A premise may be informative without
logically settling the variable named in the question. In (6.11)--(6.15),
by contrast, the premises also establish coverage of a nonpolar list.
This source-native treatment is distinct from the later PMCE proof
presentation, even where their classical semantic instances agree.

\subsection{All first-order evocation examples}

\begin{hbtheorem}{thm:p18-ep-firstorder-polar}{proved=annex:p18.ep.firstorder-polar}{Finite observations, universal claims and identity: (6.16)--(6.19)}
For each finite \(n\geq1\), observations \(Pc_1,\ldots,Pc_n\)
evoke \(?\forall xPx\); observations
\(Pc_1\land Rc_1,\ldots,Pc_n\land Rc_n\) evoke
\(?\forall x(Px\supset Rx)\). Adding \(Pc_{n+1}\) to the latter
observations evokes \(?Rc_{n+1}\). Finally \(Pc_1,Pc_2\) evoke
\(?c_1=c_2\). These are (6.16)--(6.19), over all source \(L_2\)-models.
No domain closure or unique-names assumption is present. Both the finite
observations and the distinct-symbol condition are kept in the examples.
\end{hbtheorem}

\begin{hbtheorem}{thm:p18-ep-named-evocation}{proved=annex:p18.ep.named-evocation}{Named and residual witnesses: (6.20)--(6.24)}
The premise \(\exists xPx\) evokes
\(?\{Pc,\exists x(Px\land x\ne c)\}\), source (6.20).
For any finite list \(c_1,\ldots,c_n\), \(n>1\), of distinct constant
symbols, the premise \(\exists x(Px\land N_{\vec c}(x))\) evokes both
\(?\{Pc_1,\ldots,Pc_n\}\) and \(?S(Px)\), sources (6.21),(6.23).
The premise \(\neg\forall x(N_{\vec c}(x)\supset Px)\) evokes both
\(?\{\neg Pc_1,\ldots,\neg Pc_n\}\) and \(?S(\neg Px)\),
sources (6.22),(6.24).

For the wh-question, non-entailment is proved for every constant of the
language, not only the displayed \(c_i\). Here the premise itself
guarantees a named witness; no assumption naming every domain individual
is required. The finite explicit lists are not identified with the whole
wh-answer sets.
\end{hbtheorem}

\begin{hbtheorem}{thm:p18-ep-complete-list}{proved=annex:p18.ep.complete-list}{Every candidate complete-list answer in (6.25)}
For the family in Definition~\ref{def:p18-ep-families},
\(\{H_{\vec c^{1}}\lor\cdots\lor H_{\vec c^{m}}\}\) evokes
\(?U(Px)\). A true disjunct supplies its complete-list answer
\(W_{\vec c^{j}}\). No arbitrary direct answer
\(W_{\vec t}\) is entailed by the whole premise.
The proof checks every finite candidate list \(\vec t\), including
lists containing names outside the displayed alternatives.
\end{hbtheorem}

\begin{hbtheorem}{thm:p18-ep-evocation-boundaries}{proved=annex:p18.ep.evocation-boundaries}{Why the model and list conditions matter}
The finite-observation evocations of universal statements fail if all
domain elements are stipulated to be among the observed named individuals.
Adding a unique-names premise likewise settles the distinct-name equality
question negatively. In (6.21)--(6.24), replacing \(n>1\) by a single
name either destroys the finite question's syntax or makes a wh-answer
entailed. In (6.25), retaining only one complete-extension disjunct entails
its associated complete-list answer and therefore does not evoke the
question. Ordinary arbitrary models still need not have a named witness
to a bare \(\exists xPx\); finite-name restriction is what changes that
case in (6.23), not a global naming convention.
\end{hbtheorem}

\subsection{Pure implications and their exact analytic status}

\begin{hbtheorem}{thm:p18-ep-pure-patterns}{proved=annex:p18.ep.pure-patterns}{All instances of (8.1)--(8.3),(8.7),(8.10)--(8.11)}
Negation reversal in both directions and permutation of the two direct
answers are regular pure implications. The binary question
\(?\pm|A,B|\) purely implies \(?(A\otimes B)\) for each of
\(\land,\lor,\supset,\equiv\). Each polar question \(?A\), \(?B\)
regularly and purely implies \(?\pm|A,B|\).
All formulas and question expressions must meet
Definition~\ref{def:p18-ep-notation}; no independence assumption is needed
for these implication assertions. If \(A,B\) are independent atoms,
the reverse binary-to-compound implication is not regular for any of the
four connectives. Singleton reduction and proper-subset reduction are
therefore different even for pure Boolean patterns.
\end{hbtheorem}

\begin{hbtheorem}{thm:p18-ep-general-refinement}{proved=annex:p18.ep.general-refinement}{The full formula extent of (8.4)--(8.5)}
For arbitrary formulas A,B,C such that both displayed questions are well
formed, \(?\{A,B\lor C\}\) regularly and purely implies
\(?\{A,B,C\}\). The reverse pure implication is also valid.
Both directions are analytic. Its reverse need not be regular, as the
distinct-atom case shows; regularity can hold in some dependent formula
instances. This supplies the source's general formula schemata, not only
the previously displayed distinct-atom example.
\end{hbtheorem}

\begin{hbtheorem}{thm:p18-ep-analytic-profile}{proved=annex:p18.ep.analytic-profile}{Analytic classifications with formula-instance boundaries}
The patterns (8.1)--(8.3) are analytic. The binary-to-compound pattern
(8.7) is analytic exactly for \(\otimes=\land\).
Pattern (8.10), \(?A\) to \(?\pm|A,B|\), is analytic exactly when
both \(B\) and \(\neg B\) belong to \(\mathcal H(?A)\).
Pattern (8.11) is analytic exactly when both \(A\) and \(\neg A\)
belong to \(\mathcal H(?B)\).
Thus both are non-analytic for independent distinct atoms. They can be
analytic for distinct compound-formula instances: \(A=p\land q,B=p\)
in (8.10), and \(B=p\land q,A=p\) in (8.11).
The source's non-analytic tags for these generic patterns are not a
universal non-analyticity theorem for every substitution instance.
\end{hbtheorem}

\subsection{A fixed response and finite answer partitions}

\begin{hbtheorem}{thm:p18-ep-signed-cofactors}{proved=annex:p18.ep.signed-cofactors}{Both signs and every connective in (8.12)--(8.13)}
Fix either sign of \(A\). On that one-premise basis,
\(?(A\otimes B)\) regularly implies \(?B\), for each of the four
connectives. Fix either sign of \(B\); then the same principal question
regularly implies \(?A\). These are all eight sign/connective cases
in each of (8.12),(8.13). A fixed sign can already settle the principal
compound: for instance \(\neg A\) entails \(\neg(A\land B)\).
The implication assertion does not say the auxiliary polar question is
evoked by that basis or that its responses settle different original answers.
\end{hbtheorem}

\begin{hbtheorem}{thm:p18-ep-polar-refinements}{proved=annex:p18.ep.polar-refinements}{Three-answer and covered-list patterns: (8.14)--(8.15)}
The question \(?\{\neg A,A\land B,A\land\neg B\}\) purely implies
\(?A\), with witnesses \(\{A\land B,A\land\neg B\}\) and
\(\{\neg A\}\). For every \(n>1\) and \(i\in J_n\),
\(?\{A_1,\ldots,A_n\}\) implies \(?A_i\) on \(\{D_n\}\),
with witnesses \(\{A_i\}\) and \(\{A_j:j\ne i\}\).
The latter is regular for \(n=2\). With independent atoms it is not
regular for \(n>2\); the three-answer pattern is not regular for
independent \(A,B\). These failures concern singleton witnesses, not
the validity of the displayed proper subsets.
\end{hbtheorem}

\begin{hbtheorem}{thm:p18-ep-denial}{proved=annex:p18.ep.denial}{Each negative response in (8.16)}
On basis \(\{D_n,\neg K_n\}\), the question
\(?\{A_1,\ldots,A_n\}\) implies
\(?\{\neg A_1,\ldots,\neg A_n\}\).
For target answer \(\neg A_i\), its witness is exactly the nonempty
proper set \(\{A_j:j\ne i\}\). The implication is regular if
\(n=2\), and is not in general regular for \(n>2\).
Neither the coverage premise nor the exclusion of simultaneous truth
may be omitted in general.
\end{hbtheorem}

\begin{hbtheorem}{thm:p18-ep-partition-patterns}{proved=annex:p18.ep.partition-patterns}{Every split and both responses in (8.18)--(8.21)}
For every \(1\leq i<n\), the four partition patterns in the sheet are
valid. In (8.18),(8.21), response \(B\) has witness
\(\{A_1,\ldots,A_i\}\), and \(\neg B\) has witness
\(\{A_{i+1},\ldots,A_n\}\). In (8.19),(8.20), the witnesses of
\(B,C\) are those same two blocks. Both blocks are nonempty and
proper. All four relations are regular for \(n=2\), but not generally
regular for larger independent answer lists.
In (8.21) the first premise is the single conditional
\(D^-_i\supset B\); no standalone \(D^-_i\) premise is inserted.
The coverage premise \(D_n\) remains separate.
\end{hbtheorem}

\subsection{Arbitrary first-order formulas and a wh-to-polar reduction}

\begin{hbtheorem}{thm:p18-ep-arbitrary-formulas}{proved=annex:p18.ep.arbitrary-formulas}{The full formula extent of (8.22)}
Let \(Ax,Bx,Cx\) be arbitrary monadic first-order formulas having
\(x\) as their only free variable, not merely atomic predicate instances.
For every constant \(c\), on
\(\{\forall x(Ax\equiv(Bx\land Cx))\}\), the polar question
\(?A(x/c)\) regularly implies
\(?\pm|B(x/c),C(x/c)|\).
The formulas may themselves contain Boolean connectives, equality,
bound variables and quantifiers. Substitution does not change bound
occurrences of \(x\). No global naming or unique-names condition is
needed for the statement, which concerns the interpretation of this
particular constant.
\end{hbtheorem}

\begin{hbtheorem}{thm:p18-ep-wh-polar}{proved=annex:p18.ep.wh-polar}{Every constant instance and response in (8.26)}
Let \(Ax\) be any formula of the stated kind, and let
\(c_1,\ldots,c_n\), \(n>1\), be distinct constant symbols.
For every \(i\in J_n\), on
\[
 X=\{\exists xAx,\ \forall x(Ax\supset N_{\vec c}(x))\},
\]
\(?S(Ax)\) implies \(?A(x/c_i)\). Its positive response has
singleton witness \(\{A(x/c_i)\}\); its negative response has witness
\(\{A(x/c_j):j\ne i\}\), a nonempty proper subset of the whole
wh-answer set. The implication is regular when \(n=2\), but is not
generally regular for \(n>2\), already for an atomic predicate.
The full basis, not the negative response alone, supplies the other
named witness. Neither \(\exists xAx\) nor the finite restriction may
be dropped in general.
\end{hbtheorem}

These examples do not require complete finite naming for every domain.
They require either a specific finite restriction on the property at issue
or a universal relation between the formula instances. The earlier
\(?S\)-to-\(?S\), \(?U\)-to-\(?U\), and wh-to-finite-list mappings
retain their separate answer types in
Theorem~\ref{thm:p18-ei-wh-cases}. The pure binary warrants used by a
search scenario likewise retain their actual earlier question premise;
an implication pattern is not permission to delete that premise and invoke
unrestricted transitivity.
\index{erotetic logic!evocation patterns}
\index{erotetic logic!analytic implication}
\index{question!complete-list answer}
\index{question!named witness}
